% 第12章静电场：题 12.1--12.3 独立推导（辅导书扫描页190--191）
% 公式采用 \(...\) 和 \[...\]；教材原图依赖已在题目中注明。
\section*{第12章\quad 静电场：独立解析}
\newcommand{\ke}{\frac{1}{4\pi\varepsilon_0}}

\paragraph{习题 12.1}
记中心试探电荷为 \(q_0=+1\)（单位正电荷），
中心至任一顶点距离为 \(a/\sqrt 2\)。
位于一条对角线两端的 \(2q\) 所产生的两个作用力大小相等、
方向相反，恰好抵消。
另一条对角线两端分别是 \(q\) 和 \(-4q\)。
前者对中心正电荷斥力指向 \(-4q\) 角，
后者对它的引力也指向该角，故
\[
 F=\ke\frac{q q_0}{a^2/2}
   +\ke\frac{4q q_0}{a^2/2}
 =\frac{10q q_0}{4\pi\varepsilon_0a^2}
 =\frac{5q q_0}{2\pi\varepsilon_0a^2}.
\]
如果题中的 \(q>0\)，力指向电荷 \(-4q\) 所在顶点；
若 \(q\) 被赋予其他符号，则须相应重定电荷符号与方向。

\paragraph{习题 12.2}
选取任一顶点处的负电荷。
其他两个顶点处的两个负电荷各给出大小
\(F_0=\ke q^2/r^2\) 的斥力，两力夹角为
\(60^\circ\)，合力沿该顶点的角平分线向外，
大小为 \(2F_0\cos30^\circ=\sqrt3 F_0\)。
等边三角形的重心到顶点距离是 \(r/\sqrt3\)；
中心正电荷对该负电荷的吸引力沿角平分线向内，大小
\[
 F_Q=\ke\frac{qQ}{(r/\sqrt3)^2}
 =3\ke\frac{qQ}{r^2}.
\]
平衡条件 \(F_Q=\sqrt3F_0\) 给出
\[
 Q=\frac{q}{\sqrt3}=\frac{\sqrt3}{3}q.
\]
由对称性，满足一个顶点的平衡即可使三个顶点同时平衡；
中央电荷受三个顶点电荷的合力也为零。

\paragraph{习题 12.3}
设四条等长的无质量细线长度都是 \(r\)，
四球构成图示菱形，上下对角线长 \(2r\cos\alpha\)，
左右对角线长 \(2r\sin\alpha\)，各边张力为 \(T\)
（在平衡构型中四条边的张力由对称性相等）。
对上端 \(-q\) 球写沿竖直对角线的力平衡：
另一下端 \(-q\) 球的斥力向上，
左右两个 \(Q\) 球的吸力与细线张力沿竖直向下，
即
\[
 \ke\frac{q^2}{4r^2\cos^2\alpha}
 =2\left(\ke\frac{qQ}{r^2}+T\right)\cos\alpha.
\]
对右端 \(Q\) 球沿水平对角线平衡，同理
\[
 \ke\frac{Q^2}{4r^2\sin^2\alpha}
 =2\left(\ke\frac{qQ}{r^2}+T\right)\sin\alpha.
\]
消去共同因子 \(2(\ke qQ/r^2+T)\)，得到
\[
 \frac{q^2}{\cos^3\alpha}
 =\frac{Q^2}{\sin^3\alpha}
 \quad\Longrightarrow\quad
 \boxed{\cot^3\alpha=\frac{q^2}{Q^2}}.
\]
核验：该角度关系是\textbf{静力平衡的必要条件}，
且形式上在 \(q=Q\) 时得到
\(\alpha=45^\circ\)。但若四根线只是不能承受压缩的柔性细线，
还须满足张力 \(T\ge0\)：
\[
 T=\frac{1}{4\pi\varepsilon_0r^2}
 \left(\frac{q^2}{8\cos^3\alpha}-qQ\right)\ge0.
\]
例如 \(q=Q\)、\(\alpha=45^\circ\) 时，
上式给出 \(T<0\)，因此四根\emph{仅能受拉}的细线实际上无法
维持那个对称静态构型。原习题的结论应理解为
“若存在所设平衡，其角度必满足该关系”，
而不能单凭角度关系断言任意 \(q,Q\) 均有可实现的平衡。
这属于原题模型条件需要进一步说明的审校风险。


% 习题12.4--12.23：由原题独立复算，配套辅导PDF约191--198页交叉参照。
% 所有“图示依赖”与近似假设仍列在原题/台账，未作独立专家认证。
\subsection*{习题12.4--12.23：独立解析}
\paragraph{习题 12.4}
按照题设，把两夸克近似视为相距
\(r=1.0\times10^{-15}\,\mathrm m\) 的点电荷，
\(q_u=2e/3\)、\(q_{\bar d}=e/3\)。
库仑定律给出
\[
 F=\frac{1}{4\pi\varepsilon_0}\frac{2e^2}{9r^2}
 =\frac{(8.9876\times10^9)\,2(1.6022\times10^{-19})^2}
 {9(1.0\times10^{-15})^2}
 \approx\boxed{51.3\,\mathrm N}.
\]
二者同号，静电力相斥；核内的真实强相互作用并未计入。

\paragraph{习题 12.5}
单个核子或电子的最大电荷误差按题意为
\(\delta q=10^{-21}e\)，氧原子含8个质子、8个中子、
8个电子，最坏同号偏差可叠加，故
\[
 |Q|_{\max}=24\times10^{-21}e
 \approx3.845\times10^{-39}\,\mathrm C.
\]
两个氧原子按质点近似，质量
\(m_{\rm O}\simeq16u\)、
\(u=1.66054\times10^{-27}\,\mathrm{kg}\)。
相距 \(r\) 时两种力的大小之比与 \(r\) 无关：
\[
 \frac{F_e}{F_g}
 =\frac{|Q_1Q_2|}{4\pi\varepsilon_0 G m_{\rm O}^2}
 \le\frac{(24\times10^{-21}e)^2}
 {4\pi\varepsilon_0G(16u)^2}
 \approx2.82\times10^{-6}.
\]
即使电荷为同号，静电斥力也远小于万有引力，
所以在题设误差上界和质点模型下净力仍相吸。
\textbf{注意：}这里的误差单位是 \(10^{-21}e\)，
不是 \(10^{-21}\,\mathrm C\)，误读会颠倒结论。

\paragraph{习题 12.6}
设电偶极子正、负电荷为 \(\pm q\)，
分别位于轴上 \(x=\pm l/2\)。
在轴外 \(x=r>l/2\) 的场强沿偶极矩方向，
\[
 E_x=\frac{q}{4\pi\varepsilon_0}
 \left[\frac1{(r-l/2)^2}-\frac1{(r+l/2)^2}\right]
 =\frac{q}{4\pi\varepsilon_0}
 \frac{2rl}{(r^2-l^2/4)^2}.
\]
当 \(r\gg l\)，分母近似 \(r^4\)，
所以
\[
 \boxed{E_x\simeq\frac{2ql}{4\pi\varepsilon_0r^3}
 =\frac{2p}{4\pi\varepsilon_0r^3}}.
\]
近似误差的相对量级为 \(O(l^2/r^2)\)。

\paragraph{习题 12.7}
在电偶极子远场近似中，电矩沿 \(x\) 轴，
\(\boldsymbol p=p\boldsymbol i\)，
\(\boldsymbol r=x\boldsymbol i+y\boldsymbol j\)。
偶极场的向量表达式
\[
 \boldsymbol E=\frac1{4\pi\varepsilon_0r^3}
 [3(\boldsymbol p\cdot\hat{\boldsymbol r})\hat{\boldsymbol r}
 -\boldsymbol p],\quad r^2=x^2+y^2
\]
来自正负点电荷场的一阶展开。
因 \(\boldsymbol p\cdot\hat{\boldsymbol r}=px/r\)，
逐分量整理得
\[
 \boxed{E_x=
 \frac{p(2x^2-y^2)}
 {4\pi\varepsilon_0(x^2+y^2)^{5/2}},\qquad
 E_y=\frac{3pxy}
 {4\pi\varepsilon_0(x^2+y^2)^{5/2}}}.
\]
检查：轴上 \(y=0\) 给 \(E_x=2p/(4\pi\varepsilon_0|x|^3)\)；
赤道 \(x=0\) 给 \(E_x=-p/(4\pi\varepsilon_0|y|^3)\)。

\paragraph{习题 12.8}
无限长线电荷 \(\lambda\) 在距离 \(r\) 处产生
\(|E|=|\lambda|/(2\pi\varepsilon_0r)\)。
另一根线相距 \(r=2a\)，其每单位长度受力大小
\[
 \boxed{\frac FL=\frac{\lambda^2}{4\pi\varepsilon_0a}}.
\]
两线带异号电荷，方向互相吸引。

\paragraph{习题 12.9}
把带电线段沿 \(x\) 轴取为
\(-L/2\le x\le L/2\)，观察点位于 \(x=r>L/2\)。
各元电荷产生同向轴向场，
\[
 E=\frac{\lambda}{4\pi\varepsilon_0}
 \int_{-L/2}^{L/2}\frac{dx}{(r-x)^2}
 =\frac{\lambda}{4\pi\varepsilon_0}
 \left[\frac1{r-L/2}-\frac1{r+L/2}\right].
\]
因此
\[
 \boxed{E=\frac{\lambda L}
 {4\pi\varepsilon_0[r^2-(L/2)^2]}}.
\]
正 \(\lambda\) 时沿正轴；远场 \(r\gg L\)
退化为总电荷 \(Q=\lambda L\) 的点电荷场。

\paragraph{习题 12.10}
令圆环上角度为 \(\theta\) 的小电荷元
\(dq=\lambda_0\sin\theta\,R\,d\theta\)。
中心处的电场元
\(d\boldsymbol E=
-\frac{dq}{4\pi\varepsilon_0R^2}
(\cos\theta\,\boldsymbol i+\sin\theta\,\boldsymbol j)\)。
全周积分：
\[
 E_x=-\frac{\lambda_0}{4\pi\varepsilon_0R}
 \int_0^{2\pi}\sin\theta\cos\theta\,d\theta=0,
\]
\[
 E_y=-\frac{\lambda_0}{4\pi\varepsilon_0R}
 \int_0^{2\pi}\sin^2\theta\,d\theta
 =-\frac{\lambda_0}{4\varepsilon_0R}.
\]
故 \(\lambda_0>0\) 时沿负 \(y\) 方向，
\(\boxed{\boldsymbol E=-\lambda_0\boldsymbol j/(4\varepsilon_0R)}\)。

\paragraph{习题 12.11}
以缺口的角平分线为 \(+x\) 轴。
缺口弧长 \(s=0.020\,\mathrm m\)，
圆半径 \(R=0.5\,\mathrm m\)，
残余塑料杆线密度
\(\lambda=Q/(2\pi R-s)\)。
完整带电圆环在圆心处场为零；
因此剩余圆弧的场等于缺失小圆弧
同密度电荷产生场的相反数。
设缺口张角 \(\Delta=s/R\)，
\[
 E_{\rm arc}=
 \frac{1}{4\pi\varepsilon_0}
 \frac{\lambda}{R}\int_{-\Delta/2}^{\Delta/2}
 \cos\theta\,d\theta
 =\frac{2\lambda}{4\pi\varepsilon_0R}
 \sin\frac{\Delta}{2}.
\]
代入原题数值得
\[
 \boxed{E\approx0.719\,\mathrm{N/C}},
\]
方向沿缺口角平分线\emph{指向缺口}。
若缺口长度被解释为端点直线距离，
其张角应改为 \(2\arcsin[s/(2R)]\)，
在此小角度下差异可忽略；此处明确采用弧长模型。

\paragraph{习题 12.12}
设线电荷密度为 \(\lambda\)，半圆半径 \(a\)。
对圆心 \(O\)，半圆位于左侧，
场的 \(x\) 分量为
\[
 E_{{\rm semi},x}
 =-\frac{\lambda}{4\pi\varepsilon_0a}
 \int_{\pi/2}^{3\pi/2}\cos\theta\,d\theta
 =+\frac{2\lambda}{4\pi\varepsilon_0a}.
\]
上、下两条由 \(x=0\) 向右延伸的半无限直线
的竖直分量相消，水平分量相加：
\[
 E_{{\rm lines},x}=
 -2\frac{\lambda}{4\pi\varepsilon_0}
 \int_0^\infty\frac{x\,dx}{(x^2+a^2)^{3/2}}
 =-\frac{2\lambda}{4\pi\varepsilon_0a}.
\]
因此 \(\boxed{\boldsymbol E(O)=\boldsymbol0}\)。

\paragraph{习题 12.13}
在题设半球面的\emph{平面底面}的某一点
\(P\) 作关于该底面的镜像，把均匀带电的另一半球面补上，
便得到一个均匀带电完整球面。
高斯定律与球对称性给出球面内任一点
\(\boldsymbol E_{\rm full}=\boldsymbol0\)。
在底面任一点，两半球面产生的场
\(\boldsymbol E_+\) 与 \(\boldsymbol E_-\) 的
平行于底面的分量因镜像而\emph{相等}，
垂直底面的分量则互为相反数。
因为总场为零且切向分量相等，
必有 \(2\boldsymbol E_{\parallel}=0\)。
故单个半球面在假想圆形底面上的场没有切向分量，
即\(\boxed{\boldsymbol E\perp\text{假想圆面}}\)。
本证明不能被误写成场必垂直于半球的\emph{曲面}。

\paragraph{习题 12.14}
(1) 点电荷位于立方体中心，按高斯定律，
穿过六个面的总电通量 \(q/\varepsilon_0\)，
由对称性各面相同：
\[
 \boxed{\Phi_{\rm face}=\frac{q}{6\varepsilon_0}}.
\]
(2) 把8个相同立方体沿点电荷所在的顶点拼成大立方体，
该点电荷位于大立方体中心。
大立方体每个外大面的通量 \(q/(6\varepsilon_0)\)，
每个大面含4个小面，故原小立方体
\emph{与点电荷不相邻的3个外面}各通量
\[
 \boxed{\Phi_{\rm outer}=\frac{q}{24\varepsilon_0}}.
\]
\emph{经过点电荷所在顶点的3个面}，
电场在面内切向（除奇点自身），故其面通量为零。
奇点不能直接用局部面元的普通积分处理；
拼合构造避免了其边界歧义。

\paragraph{习题 12.15}
由导体地表外侧法向场的跳跃条件
\(\sigma=\varepsilon_0E_n\)，
外法线取向上，而晴天场向下，
因此地表面电荷密度为负，
\(\sigma=-\varepsilon_0(120\,\mathrm{V/m})\)。
每平方米过剩电子数
\[
 \boxed{N=\frac{|\sigma|}{e}
 =\frac{\varepsilon_0(120)}{1.60218\times10^{-19}}
 \approx6.63\times10^9\ \mathrm{m^{-2}}}.
\]
所指的是负电荷过剩而非正电荷过剩。

\paragraph{习题 12.16}
把竖直向上取为 \(+z\)，
\(E_z(100\,\mathrm m)=-150\,\mathrm{N/C}\)，
\(E_z(300\,\mathrm m)=-100\,\mathrm{N/C}\)。
对面积 \(A\)、高度 \(\Delta z=200\,\mathrm m\)
的高斯柱，侧面通量为零，
\[
 \rho_{\rm avg}
 =\varepsilon_0\frac{E_z(300)-E_z(100)}{200}
 =\varepsilon_0\frac{50}{200}
 \approx2.214\times10^{-12}\,\mathrm{C/m^3}.
\]
故是\emph{正}体电荷密度，即每立方米
\[
 \boxed{n_{\rm deficit}=\rho_{\rm avg}/e
 \approx1.38\times10^7}
\]
个电子的净缺少量。
单位检查：\(\varepsilon_0 E/\ell\) 是 \(\mathrm{C/m^3}\)。

\paragraph{习题 12.17}
对半径 \(r\)、长度 \(\ell\) 的同轴高斯柱面，
圆柱面通量为 \(E(r)2\pi r\ell\)。
若 \(r<a\)，包围净电荷为零；
若 \(r>a\)，包围电荷为
\(\sigma(2\pi a\ell)\)。
因而
\[
 \boxed{E_r(r)=
 \begin{cases}
 0,&0\le r<a,\\[2pt]
 \displaystyle\frac{\sigma a}{\varepsilon_0r},&r>a.
 \end{cases}}
\]
方向正 \(\sigma\) 时径向向外；在 \(r=a\)
电场法向分量跨越表面有跳跃
\(E_r(a^+)-E_r(a^-)=\sigma/\varepsilon_0\)。
\(E-r\) 曲线为内部零、表面跃变、
外部按 \(1/r\) 衰减。

\paragraph{习题 12.18}
两层同轴圆筒的线电荷密度分别为
\(+\lambda,-\lambda\)，取 \(R_1<R_2\)。
高斯柱面每单位长度所包围的电荷依次为
\(0,\lambda,\lambda-\lambda=0\)，故
\[
 \boxed{E_r(r)=
 \begin{cases}
 0,&r<R_1,\\
 \displaystyle\frac{\lambda}{2\pi\varepsilon_0r},
 &R_1<r<R_2,\\
 0,&r>R_2.
 \end{cases}}
\]
在筒面上取两侧极限，不能把薄层表面的跃变
误作内部连续值。

\paragraph{习题 12.19}
设第一平面位于 \(x=0\)，第二平面位于
\(x=d>0\)，法向 \(+\boldsymbol i\) 从第一面指向第二面。
无限带电平面两侧的场大小
\(|\sigma|/(2\varepsilon_0)\)，
符号由其朝向决定。
叠加后
\[
 \boxed{E_x=
 \begin{cases}
 -(\sigma_1+\sigma_2)/(2\varepsilon_0),&x<0,\\
 (\sigma_1-\sigma_2)/(2\varepsilon_0),&0<x<d,\\
 (\sigma_1+\sigma_2)/(2\varepsilon_0),&x>d.
 \end{cases}}
\]
代入 \(\sigma_1=4\times10^{-11}\),
\(\sigma_2=-2\times10^{-11}\,\mathrm{C/m^2}\)：
\[
 E_x\simeq
 \begin{cases}
 -1.13\,\mathrm{N/C},&x<0,\\
 +3.39\,\mathrm{N/C},&0<x<d,\\
 +1.13\,\mathrm{N/C},&x>d.
 \end{cases}
\]
这里第一面右侧正向，第二面左侧因带负电而指向第二面。

\paragraph{习题 12.20}
厚壁中心为 \(d=0\)，表面在 \(d=\pm D/2\)。
在距中面 \(|d|<D/2\) 处取对称的高斯柱，
通量 \(2E(d)A\)，包围电荷
\(2\rho dA\)（取 \(d>0\)），
所以 \(E(d)=\rho d/\varepsilon_0\)。
超出厚壁后包围总电荷等效面密度
\(\rho D\)：
\[
 \boxed{E_d(d)=
 \begin{cases}
 \rho d/\varepsilon_0,&|d|\le D/2,\\
 \rho D\,\operatorname{sgn}(d)/(2\varepsilon_0),
 &|d|>D/2.
 \end{cases}}
\]
带正电时在上下两侧向外；
\(E-d\) 曲线在厚壁内部是过原点的直线，
壁外为正、负常值平台。

\paragraph{习题 12.21}
设开孔的无限平面位于 \(z=0\)，
孔半径 \(R\)，面电荷密度 \(\sigma_0\)。
用完整均匀无限平面\emph{减去}
中心半径为 \(R\) 的带电圆盘。
对 \(z>0\)，圆盘在轴上的场为
\[
 E_{\rm disk}(z)=\frac{\sigma_0}{2\varepsilon_0}
 \left(1-\frac{z}{\sqrt{z^2+R^2}}\right),
\]
无限平面场为 \(\sigma_0/(2\varepsilon_0)\)，
故孔轴上
\[
 \boxed{E_z(z)=\frac{\sigma_0z}
 {2\varepsilon_0\sqrt{z^2+R^2}}}.
\]
对 \(z<0\) 通过奇对称性延拓，同一表达式仍成立。
检查：\(E_z(0)=0\)，远轴 \(|z|\gg R\) 时
趋向有孔无限面原来的 \(\pm\sigma_0/(2\varepsilon_0)\)。

\paragraph{习题 12.22}
原均匀大球球心为 \(O\)，空腔球心为 \(O'\)，
其相对位移向量 \(\boldsymbol a=\overrightarrow{OO'}\)。
把“挖空”看成叠加：
原密度为 \(+\rho\) 的完整大球，及以
\(O'\) 为中心密度为 \(-\rho\) 的小球。
任何空腔内位置 \(P\) 的位矢分别为
\(\boldsymbol R=\overrightarrow{OP}\) 与
\(\boldsymbol r=\overrightarrow{O'P}\)，
且 \(\boldsymbol R=\boldsymbol a+\boldsymbol r\)。
均匀带电实心球内部电场为
\(\rho\boldsymbol r/(3\varepsilon_0)\)，
叠加得
\[
 \boldsymbol E_{\rm cavity}
 =\frac{\rho}{3\varepsilon_0}
 (\boldsymbol R-\boldsymbol r)
 =\boxed{\frac{\rho\,\boldsymbol a}{3\varepsilon_0}}.
\]
与 \(P\) 的位置无关，因此空腔内均匀。
若 \(\rho<0\)，方向与 \(\boldsymbol a\) 相反。

\paragraph{习题 12.23}
设电子云密度 \(\rho_e=-C e^{-2r/a_0}\)，
总电子电荷必须是 \(-e\)：
\[
 -e=-4\pi C\int_0^\infty r^2e^{-2r/a_0}\,dr
 =-\pi C a_0^3.
\]
故 \(\boxed{C=e/(\pi a_0^3)}\)。
半径 \(r\) 内的电子云所包围电荷为
\[
 Q_e(r)=-e\left[1-e^{-2r/a_0}
 \left(1+\frac{2r}{a_0}+\frac{2r^2}{a_0^2}\right)\right].
\]
再计入中心的质子电荷 \(+e\)，
球内\emph{净}电荷是
\[
 Q_{\rm net}(r)=e\,e^{-2r/a_0}
 \left(1+\frac{2r}{a_0}+\frac{2r^2}{a_0^2}\right).
\]
在 \(r=a_0=0.53\times10^{-10}\,\mathrm m\) 时
\[
 \boxed{Q_{\rm net}(a_0)=5e^{-2}e
 \approx0.67668e\approx1.084\times10^{-19}\,\mathrm C}.
\]
按球对称高斯定律，
\[
 \boxed{E(a_0)=\frac{5e^{-2}e}
 {4\pi\varepsilon_0a_0^2}
 \approx3.47\times10^{11}\,\mathrm{N/C}},
\]
方向径向向外。
若只计算电子云的场而漏掉质子，便不符合题目
“半径内净电荷”的中性原子模型。
